\section{Formative Study: Understanding and Scoping Skill Repair}
\label{sec:formative}

Repair episodes expose a recurring asymmetry in agent use: workflow owners can
recognize an unacceptable execution while remaining uncertain how the skill
produced it or how future cases should change. We use \emph{workflow owner} to
describe a person who defines the goals of a recurring task, judges the agent's
behavior, or bears responsibility for its consequences. We conducted an
artifact-grounded formative study to examine how workflow owners interpret
skill behavior, construct the scope of a repair, and divide repair work with
AI. We asked:

\begin{itemize}[leftmargin=*,itemsep=2pt,topsep=3pt]
    \item \textbf{FQ1:} How do workflow owners relate an observed breakdown to
    the instructions and interactions within a reusable skill?

    \item \textbf{FQ2:} How do additional cases and execution evidence reshape
    what workflow owners believe should change, remain stable, or remain
    unresolved?

    \item \textbf{FQ3:} How do workflow owners distribute responsibility
    between themselves and AI across evidence gathering, repair
    implementation, and release?
\end{itemize}

\subsection{Participants, Procedure, and Analysis}
\label{sec:formative-method}

We recruited \StudyN{} workflow owners who had used an LLM assistant or agent
for the same recurring task at least five times and had encountered an
unacceptable execution within the previous six months. Eligibility centered on
responsibility for the workflow's goals or outcomes; participants did not need
to author its reusable instructions. Maximum-variation sampling covered
different levels of familiarity with inspecting or modifying instructions and
different consequences of incorrect behavior. Some participants primarily
evaluated domain behavior and relied on a teammate or AI to change the
configuration, while others regularly maintained system prompts, templates,
routines, or skills. Before the session, participants provided a redacted
failure execution, a typical or successful execution from the same workflow,
and any available instruction, configuration, or revision artifact. Sessions
lasted approximately 80 minutes; participants received \StudyCompensation{},
and the protocol was \StudyEthics{}. Appendix~\ref{app:participants} reports
recruitment and sample details.

Each session included three activities. First, participants used their own
materials in an \emph{artifact-grounded breakdown reconstruction} (25 minutes)
to reconstruct the execution, their expectations, the evidence they observed,
their inferences, and any changes they considered. The interviewer used a
neutral event timeline and introduced no intent--skill--workflow model or
candidate instruction. Second, participants completed an \emph{open-ended
Wizard-of-Oz repair probe} (35 minutes) built around a common travel-recovery
incident. After seeing the request and observable execution, participants
received the reusable skill and available tools. They could request another
context, repeated executions, a rule-level what-if, tool-state evidence, or a
comparison of candidate revisions. A facilitator retrieved each response from
a fixed, precomputed execution bank. Every participant then received a
counterbalanced evidence set with a preservation case, a difficult boundary
case, repeated runs, and an instruction-level what-if. Participants chose to
write a revision, request an AI draft, specify the desired behavior, or defer
the edit, followed by a release, revision, further-evidence, or defer decision.
Third, a \emph{reflection and delegation interview} (15 minutes) compared the
probe with participants' own practices and examined preferred representations,
implementation support, and decision authority. Appendices
\ref{app:procedure}--\ref{app:instruments} provide the complete materials.

We collected submitted artifacts, transcripts, event timelines, evidence
requests, successive explanations and repair proposals, behavioral
predictions, implementation choices, and final decisions. Our analysis had two
stages. Inductive event coding first characterized observed symptoms, expected
behavior, candidate explanations, evidence sources, uncertainty, proposed
actions, scope revisions, and decision rationales. For FQ2, we built
within-participant process traces linking evidence to changes in interpretation,
repair scope, confidence, and action. For FQ3, we compared participants' choices
in the probe with the responsibilities they described in interviews. We then
conducted an abductive cross-case synthesis of how participants connected
intended behavior, reusable instructions, and enacted workflows. Two
researchers independently examined an initial maximally varied subset, refined
the coding and event schema, and searched the full corpus for disconfirming
cases. We analyzed real incidents and standardized-probe behavior separately
before comparing them. Descriptive counts characterize patterns within our
sample. Appendix~\ref{app:analysis} details the analysis and claim-audit
process.

\subsection{Findings}
\label{sec:formative-findings}

% The following titles are provisional analytic propositions. Retain, soften,
% or replace them only after the full analysis and negative-case review.

\subsubsection{F1. Skill text alone left instruction effects unclear}
\label{sec:f1-locus}

Across participants' real incidents, \DataNeeded{X of N / NUMBER OF INCIDENTS}
workflow owners could readily explain why an observed outcome was unacceptable,
yet lacked a stable account of what should change in the skill. Their accounts
included an underspecified policy, one instruction, an interaction among
otherwise reasonable instructions, context-dependent activation, and variable
execution. Tool results also mattered because they supplied values consumed by
the instructions. Participants moved among these explanations, with no direct
progression from a failed output to a faulty line of text. Their accounts varied in
evidential status: traces exposed some actions or tool results, whereas claims
about rule priority or model interpretation were often retrospective
inferences. \DataNeeded{REPORT HOW OFTEN PARTICIPANTS REVISED THE PRESUMED
RELEVANT INSTRUCTIONS OR INTERACTIONS DURING THE REAL-INCIDENT RECONSTRUCTION.}

Access to the reusable artifact changed, but did not necessarily resolve, this
uncertainty. Participants with less implementation experience tended to reason
from desired outcomes, examples, and operational consequences; experienced
maintainers more readily named instructions and interactions. However,
\DataNeeded{REPORT WHICH DIFFICULTIES PERSISTED AMONG EXPERIENCED MAINTAINERS,
SUCH AS PRIORITY, CONTEXT-DEPENDENT ACTIVATION, OR STOCHASTICITY}. The
standardized probe
provided a common illustration: after seeing the same incident and later the
same skill, participants \DataNeeded{REPORT THE RANGE OF INSTRUCTIONS OR
INTERACTIONS THEY CONSIDERED AND HOW MANY CHANGED THEIR EXPLANATION}. A
representative participant distinguished
behavioral judgment from implementation certainty: \emph{``\DataNeeded{VERIFIED
QUOTE ABOUT KNOWING THE BEHAVIOR WAS WRONG BUT NOT KNOWING WHICH INSTRUCTION OR
INTERACTION SHOULD CHANGE}.''}

Negative cases bounded this pattern. \DataNeeded{REPORT INCIDENTS IN WHICH THE
PARTICIPANT RAPIDLY IDENTIFIED THE RELEVANT INSTRUCTION OR INTERACTION}. Taken
together, the accounts separated workflow accountability from artifact
expertise. The skill provided an editable control surface, while execution
evidence helped participants connect its instructions to observable behavior.

\subsubsection{F2. Repair scope was constructed comparatively across cases}
\label{sec:f2-scope}

Participants initially expressed repair intent in several forms: an objection
to the visible outcome, a preferred decision for the current case, a candidate
textual edit, or a broader policy. \DataNeeded{REPORT THE DISTRIBUTION OF THESE
INITIAL FORMS AND HOW OFTEN INITIAL PROPOSALS OMITTED AN APPLICABILITY BOUNDARY,
PRESERVATION COMMITMENT, OR UNRESOLVED TRADE-OFF}. These first proposals often
answered ``what should have happened here'' more clearly than ``which future
cases should follow the same rule.'' The real incidents showed the practical
consequence of this gap: \DataNeeded{REPORT HOW MANY PARTICIPANTS DESCRIBED A
PREVIOUS LOCAL FIX THAT LATER CHANGED ANOTHER DESIRED BEHAVIOR, DELAYED RELEASE,
OR LED TO ROLLBACK/WORKAROUND}.

The open evidence-request period let us distinguish what participants
anticipated from what the study subsequently introduced. Participants
spontaneously requested \DataNeeded{REPORT REQUEST TYPES AND FREQUENCIES:
CONTRASTING CONTEXTS, REPEATED RUNS, TOOL STATE, WHAT-IFS, PATCH COMPARISONS,
RAW TRACES}. These requests exposed different uncertainties: some participants
sought another case to discover the boundary of a rule, while others wanted
repetition to determine whether the observed behavior was stable. The common
evidence set then provided comparable pressure on every proposal. A routine
case could reveal an established behavior that a broad repair would disturb; a
boundary case could surface a trade-off for which the participant had not yet
formed a policy; repeated runs could weaken confidence in a single execution;
and fixed tool-state evidence could clarify which skill conditions had been
satisfied.

Our process traces captured whether evidence changed more than confidence.
After self-requested and common evidence, \DataNeeded{X PARTICIPANTS NARROWED,
Y BROADENED, Z REFRAMED, W DEFERRED, AND V ABANDONED THEIR INITIAL REPAIR}.
Preservation commitments first appeared \DataNeeded{WHEN/AFTER WHICH EVIDENCE},
and unresolved regions appeared \DataNeeded{WHEN/AFTER WHICH EVIDENCE}.
Participants sometimes revised their interpretation of the original workflow
at the same time as the proposed repair. As one participant explained,
\emph{``\DataNeeded{VERIFIED QUOTE SHOWING THAT ANOTHER CASE CHANGED WHAT THE
RULE SHOULD MEAN, NOT ONLY CONFIDENCE IN AN EXISTING EDIT}.''} We also observed
\DataNeeded{NEGATIVE CASES IN WHICH ADDITIONAL EVIDENCE ONLY CONFIRMED A
PREVIOUSLY ARTICULATED POLICY}.

Contrastive executions served two roles in these transitions. They tested
existing commitments and helped participants construct the relative scope of
change: where behavior should differ, where it should remain stable, and where
the desired policy remained unsettled. Participants could then narrow, defer,
or abandon a proposed skill edit when its behavioral scope lacked support.

\subsubsection{F3. Participants separated evidence and implementation work from policy and release authority}
\label{sec:f3-delegation}

Participants distributed repair work across three forms of responsibility.
\emph{Evidence work} included generating cases, repeating executions,
conducting what-if comparisons, and organizing results. \emph{Implementation
work} included drafting or applying textual revisions. \emph{Policy and release
work} included judging acceptable behavior, resolving trade-offs, deciding
whether evidence was sufficient, and authorizing reuse. In the probe,
participants actually delegated \DataNeeded{REPORT EVIDENCE TASKS AND COUNTS}
and selected \DataNeeded{COUNTS FOR MANUAL EDITING, AI DRAFTING,
BEHAVIOR-ONLY SPECIFICATION, DEFERRAL, OR NO EDIT}. These enacted choices did
not always match their initial statements about preferred control:
\DataNeeded{REPORT ANY SYSTEMATIC MISMATCH BETWEEN WHAT PARTICIPANTS SAID AND
WHAT THEY DID}.

Artifact familiarity shaped the level at which participants wanted to engage,
but it did not map cleanly onto a single allocation of authority.
\DataNeeded{REPORT HOW LESS-EXPERIENCED WORKFLOW OWNERS AND EXPERIENCED
MAINTAINERS DIFFERED IN THEIR DESIRE FOR TEXTUAL DETAIL, RAW EXECUTIONS, OR
MANUAL EDITING}. Perceived consequence also mattered: \DataNeeded{REPORT WHEN
HIGHER CONSEQUENCE PRODUCED MORE CONSERVATIVE DELEGATION OR STRONGER
TRACEABILITY DEMANDS}. At the end of the probe, participants chose
\DataNeeded{COUNTS FOR RELEASE, REVISE, MORE EVIDENCE, AND DEFERRAL}
and justified these choices in terms of \DataNeeded{DOMINANT DECISION
CRITERIA}. A participant located their desired authority precisely:
\emph{``\DataNeeded{VERIFIED QUOTE DISTINGUISHING WILLINGNESS TO DELEGATE
EXPERIMENTS OR PATCH DRAFTING FROM AUTHORITY OVER BEHAVIOR OR RELEASE}.''}

The accounts varied. \DataNeeded{REPORT PARTICIPANTS WHO INSISTED
ON DIRECT TEXTUAL AUTHORSHIP, THOSE WILLING TO DELEGATE MOST IMPLEMENTATION,
AND ANY WILLING TO DELEGATE POLICY OR RELEASE}. Participants negotiated control
across behavioral commitments, textual implementation, access to evidence, and
the release decision. This division of labor motivates a workflow in which AI
can perform evidence and implementation work while users retain access to the
artifact and authority over whether a behavioral change is ready for reuse.

\paragraph{Synthesis: repair as intent--skill--execution reconciliation.}
Together, the findings characterize repair as an iterative process
(Figure~\ref{fig:repair-process}). Workflow owners interpret the current skill,
construct a provisional scope across cases, inspect relevant instructions,
implement a candidate revision, and review its behavioral effects. New evidence
may revise an earlier assumption about the instructions or the policy boundary.
The skill thus becomes an editable intermediary between user intent and enacted
workflow behavior.

\begin{figure}[t]
    \centering
    \begin{tikzpicture}[
        node distance=4.5mm,
        every node/.style={font=\scriptsize,align=center},
        box/.style={draw,rounded corners=2pt,fill=draftblue,
                   text width=0.82\columnwidth,minimum height=7mm,inner sep=3pt},
        arrow/.style={-{Latex[length=1.6mm]},line width=0.45pt}
    ]
        \node[box] (a) {Interpret current skill behavior\\
        \emph{instructions, interactions, and executions}};
        \node[box,below=of a] (b) {Construct a provisional repair scope\\
        \emph{change, preserve, and unresolved commitments}};
        \node[box,below=of b] (c) {Probe instructions and implement a candidate revision};
        \node[box,below=of c] (d) {Review the observed behavioral change};
        \node[box,below=of d] (e) {Release, revise, gather evidence, or defer};
        \draw[arrow] (a) -- (b);
        \draw[arrow] (b) -- (c);
        \draw[arrow] (c) -- (d);
        \draw[arrow] (d) -- (e);
        \draw[arrow,dashed] (d.east) -- ++(4mm,0) |- (b.east);
        \draw[arrow,dashed] (c.west) -- ++(-4mm,0) |- (a.west);
    \end{tikzpicture}
    \caption{Provisional process model to retain only if supported by the final
    analysis. Evidence can revise both the presumed instruction relationship and
    the intended scope of repair.}
    \Description{A vertical five-stage process: interpret current skill behavior, construct repair scope, probe instructions and implement a candidate, review behavioral change, and choose release, revision, more evidence, or deferral. Dashed arrows return from later stages to earlier interpretation and scoping.}
    \label{fig:repair-process}
\end{figure}

\subsection{Design Goals for Behavioral Patch Review}
\label{sec:formative-dgs}

We synthesized three design goals for supporting skill repair when runtime
behavior and the desired policy remain only partially understood.

\paragraph{DG1: Connect skill text to execution evidence.}
The interface should help users relate skill instructions and edits to behavior
across repeated executions. Task context and tool results provide controlled
inputs for these comparisons, and each summary should remain linked to its
executed intervention and raw runs [F1].

\paragraph{DG2: Make the repair scope explicit and revisable.}
The interface should help workflow owners articulate where a repair should
change behavior, preserve established behavior, or leave a policy unresolved.
These commitments should be grounded in concrete cases and remain editable as
new executions refine the owner's interpretation. The resulting scope serves
as both an evolving specification and the reference for reviewing a candidate
[F1, F2].

\paragraph{DG3: Support evidence-grounded review while preserving release authority.}
AI may generate cases, run probes, organize comparisons, and draft revisions.
Workflow owners should be able to see how their behavioral scope guides the
textual implementation and how the complete revision changes observable
decisions, tool actions, and outcomes. The workflow should culminate in a choice
to release, revise, gather more evidence, or defer [F2, F3].

Together, these goals motivate \emph{behavioral patch review}: a shared workflow
that connects user judgments, reusable skill instructions, execution evidence,
and enacted behavior. The next section describes how \systemname{}
operationalizes this process.
